\section{Contextualização e Cenário}

\begin{frame}{O Sistema de Fomento ao Microempreendedor (SFM)}
  \begin{itemize}
    \item \textbf{O que é:} Plataforma governamental do Estado da Paraíba em parceria com o ARIA (CI/UFPB).
    \item \textbf{Objetivo:} Intermediar a contratação de serviços de pronto pagamento entre órgãos públicos e microempreendedores (MEIs).
    \item \textbf{Dinâmica principal:} 
    \begin{itemize}
        \item Órgãos cadastram \textbf{oportunidades} (demandas).
        \item MEIs submetem \textbf{propostas} de execução.
        \item O sistema gere a seleção, aprovação e homologação.
    \end{itemize}
  \end{itemize}
\end{frame}

\begin{frame}{Dinâmica Operacional da Equipe}
  \begin{itemize}
    \item \textbf{Estrutura:} $\approx$ 21 colaboradores (em sua maioria discentes) e 2 professores coordenadores em regime de trabalho remoto.
    \item \textbf{Metodologia Ágil:}
    \begin{itemize}
        \item Reuniões gerais semanais com atas de acompanhamento.
        \item Alinhamentos periódicos por subequipes (design, \textit{back-end}, \textit{front-end}, infraestrutura).
        \item Gestão de tarefas via \textit{backlog} (Kanban/GitHub).
    \end{itemize}
    \item \textbf{Responsabilidades dos Testes:}
    \begin{itemize}
        \item A equipe de qualidade era responsável pela elaboração e execução dos testes de sistema.
        \item Testes de unidade e de integração eram de responsabilidade das equipes que implementavam as funcionalidades.
    \end{itemize}
  \end{itemize}
\end{frame}

\begin{frame}{O Fluxo de Desenvolvimento Empírico}
  \textbf{Como uma nova funcionalidade era concebida:}
  \begin{enumerate}
    \item \textbf{Ideação:} Debate conciso na reunião geral semanal.
    \item \textbf{Prototipação:} A equipe de \textit{design} elaborava as interfaces visuais.
    \item \textbf{Implementação:} \textit{Front} e \textit{Back-end} codificavam com base exclusivamente nos protótipos.
    \item \textbf{Testes (QA):} Casos de teste estruturados \textbf{apenas após} a codificação, baseados na interface gráfica e em dedução empírica.
  \end{enumerate}
  

  \vspace{1em}
  Frequentemente, o primeiro passo era ignorado, sendo de responsabilidade única da equipe de \textit{design} a elaboração da funcionalidade.
\end{frame}